% ====================================================================
% FF2-2 (ch1 §4.2 배치 제안; §13.5 도핑 절에서 재참조) — 히스 gap 문턱 곡선
%   좌표 = eq:dUhys 를 T=298.15 K 에서 python 으로 점별 수치 평가:
%   gap(Ω) = (2/F)[Ω u − 2RT artanh u], u=√(1−2RT/Ω) (Ω>2RT), 그 외 0.
%   검산: gap(4RT)=54.76 mV = fig:hysloop 의 2.131 RT/F 와 일치.
%   양수 점근 (8RT/3F)u³, 상수-대입 오답 −(4RT/3F)u³(★Taylor 함정), staging 초기값 4점,
%   도핑 극한(eq:lco-dope) 주석 포함. 사용 라이브러리: ch1 preamble 범위 내(arrows.meta).
% ====================================================================
\begin{figure}[t]
\centering
\begin{tikzpicture}[x=1.5cm,y=0.030cm]
% single-phase region (gap identically 0)
\fill[black!7] (0,-19) rectangle (2,120);
\node[font=\scriptsize,align=center] at (1.0,63) {단상 쪽\\$\Omega_j\le2RT$:\\gap $\equiv0$};
\draw[densely dotted] (2,-19) -- (2,120);
\node[font=\scriptsize,anchor=south,rotate=90] at (2.0,72) {$\Omega=2RT$ ($T{=}T_{c,j}$)};
% axes
\draw[-{Latex[length=1.6mm]}] (0,-19) -- (0,126) node[above,font=\scriptsize] {$\Delta U_j^\hys$ [mV]};
\draw[-{Latex[length=1.6mm]}] (0,0) -- (5.95,0) node[below,font=\scriptsize] {$\Omega_j/RT$};
\foreach \xx in {1,2,3,4,5} {\draw (\xx,1.7) -- (\xx,-1.7) node[below,font=\scriptsize] {\xx};}
\foreach \yy in {25,50,75,100} {\draw (0.013,\yy) -- (-0.013,\yy) node[left,font=\scriptsize] {\yy};}
% right axis in RT/F units
\draw (5.75,0) -- (5.75,120);
\foreach \yy/\lab in {25.69/1,51.39/2,77.08/3,102.77/4}
  {\draw (5.737,\yy) -- (5.763,\yy) node[right,font=\scriptsize] {\lab};}
\node[above,font=\scriptsize] at (5.75,120) {[$RT/F$]};
% exact gap curve (zero below threshold, then closed form)
\draw[very thick] (0,0) -- (2,0);
\draw[very thick] plot[smooth] coordinates {(2.00,0.00) (2.05,0.27) (2.10,0.75) (2.15,1.38) (2.20,2.10) (2.25,2.92) (2.30,3.81) (2.35,4.77) (2.40,5.80) (2.45,6.87) (2.50,8.00) (2.55,9.17) (2.60,10.38) (2.65,11.63) (2.70,12.93) (2.75,14.25) (2.80,15.61) (2.85,17.00) (2.90,18.41) (2.95,19.86) (3.00,21.33) (3.10,24.34) (3.20,27.45) (3.30,30.64) (3.40,33.90) (3.50,37.23) (3.60,40.62) (3.70,44.08) (3.80,47.59) (3.90,51.15) (4.00,54.76) (4.10,58.42) (4.20,62.11) (4.30,65.85) (4.40,69.63) (4.50,73.44) (4.60,77.29) (4.70,81.17) (4.80,85.08) (4.90,89.02) (5.00,92.98) (5.10,96.98) (5.20,101.00) (5.30,105.04) (5.40,109.10) (5.50,113.19) (5.60,117.30)};
% correct cubic onset asymptote (both series expanded together)
\draw[densely dashed] plot[smooth] coordinates {(2.00,0.00) (2.06,0.34) (2.12,0.92) (2.18,1.63) (2.24,2.40) (2.30,3.23) (2.36,4.08) (2.42,4.95) (2.48,5.83) (2.54,6.72) (2.60,7.60) (2.66,8.47) (2.72,9.33) (2.78,10.18) (2.84,11.02) (2.90,11.85) (2.96,12.65) (3.02,13.45) (3.08,14.23) (3.14,14.99) (3.20,15.73)};
\node[font=\scriptsize,anchor=west] at (3.24,14.5) {문턱 점근 $\tfrac{8RT}{3F}u_j^3>0$};
% Taylor trap: constant substitution gives wrong negative gap
\draw[densely dotted,thick] plot[smooth] coordinates {(2.00,0.00) (2.12,-0.46) (2.24,-1.20) (2.36,-2.04) (2.48,-2.92) (2.60,-3.80) (2.72,-4.67) (2.84,-5.51) (2.96,-6.33) (3.08,-7.11) (3.20,-7.87)};
\node[font=\scriptsize,anchor=west] at (2.55,-16) {★함정: 상수 대입 오답 $-\tfrac{4RT}{3F}u_j^3<0$ (비물리)};
\draw[-{Latex[length=1.2mm]},black!60] (2.9,-13) -- (3.02,-7.6);
% staging initial values on the exact curve
\fill (2.420,5.80) circle (1.3pt);
\node[font=\tiny,anchor=south east,inner sep=1.2pt] at (2.42,7.6) {$4{\to}3$ (6.2)};
\fill (3.227,28.31) circle (1.3pt);
\node[font=\tiny,anchor=south east,inner sep=1.2pt] at (3.23,30.1) {$3{\to}2\mathrm L$ (28.3)};
\fill (4.034,56.00) circle (1.3pt);
\node[font=\tiny,anchor=south east,inner sep=1.2pt] at (4.03,57.8) {$2\mathrm L{\to}2$ (56.0)};
\fill (5.244,102.78) circle (1.3pt);
\node[font=\tiny,anchor=south east,inner sep=1.2pt] at (5.24,104.6) {$2{\to}1$ (102.8)};
% cross-check with fig:hysloop at Omega=4RT
\draw[densely dotted,black!60] (4.0,7.5) -- (4.0,54.76) -- (0,54.76);
\node[font=\tiny,anchor=north west,align=left] at (4.05,44) {$(4RT,\,54.8$ mV$)$\\$=2.131\,RT/F$};
% doping arrow toward threshold (just above the x-axis, pointing left)
\draw[-{Latex[length=1.6mm]},black!60] (3.35,3.5) -- (2.30,3.5);
\node[font=\scriptsize,anchor=west,black!60] at (3.42,5.0) {도핑 $\Omega^\mathrm{cat,dop}\!\to\!2RT^{+}$: gap$\,\to0$};
\end{tikzpicture}
\caption{히스테리시스 분기 gap 의 문턱 곡선 $\Delta U_j^\hys(\Omega_j)$ --- 좌표는
식~\eqref{eq:dUhys} 그대로의 수치 평가다($T{=}298.15$ K). $\Omega_j\le2RT$(음영)에서 gap 은 정확히
$0$ 이고, 문턱 직후에는 두 급수를 함께 전개한 $\tfrac{8RT}{3F}u_j^3$(파선)으로 \emph{양수로}
소멸한다($u_j^3\propto(T_{c,j}-T)^{3/2}$ --- 임계온도에서 매끄러운 소멸). 점선은 첫째 항에
$\Omega_j\!\approx\!2RT$ 를 상수 대입만 한 오답 $-\tfrac{4RT}{3F}u_j^3$ 으로, 부호가 뒤집힌 비물리
gap 을 준다(본문 ★Taylor 전개의 함정). ● 는 표~\ref{tab:staging} 초기값 $\Omega_j$ 가 갖는 spinodal
상한(6.2/28.3/56.0/102.8 mV)이고, $\Omega{=}4RT$ 의 $54.8$ mV $=2.131\,RT/F$ 는
그림~\ref{fig:hysloop} 와 같은 값이다(교차 검산). 실측 분기는 $\gamma_j\in[0,1]$ 로 이 상한
안쪽이며(식~\eqref{eq:Ubranch}), 도핑은 $\Omega^\mathrm{cat,dop}\to2RT^{+}$ 로 같은 곡선 위에서
gap 을 닫는다(식~\eqref{eq:lco-dope}).}
\label{fig:cand-ff2-2}
\end{figure}
